% SGA 3, Expose VII_B, English reconstruction -- Loop 1.
% Scope: Theorem 4.4, Remark 4.4.A, and Corollary 4.4.1 in full
% on authority local pages 79--80; stops immediately before Section 4.4.2.
% Combined-reader pages: 592--593. Re-edition printed pages: 581--582.
% Original-edition marginal locators in scope: 549--550.
% Editorial notes in scope: 139--140.
% Authority: Polo--Gille Exp7B-13oct24.pdf,
% SHA-256 D303690E0195654BBE7C3DAA5D89EA23E24CD709A3575217234AC19AE8836694.
% Jacob Reinhold's CC BY 4.0 Markdown was consulted as comparison lineage.
% No genuine diagram occurs in this component; all mathematics is native LaTeX.

\subsection*{\SGAInlineRefTarget{sga3:viib:subsection:4-4}4.4 定理}
\phantomsection
\label{sga3:VIIB:4.4}

\noindent\textbf{定理 4.4.}
{\itshape 设 \(k\) 为特征 \(p>0\) 的伪紧环，\(G\) 为形式 \(k\)-群。下列断言等价：
\begin{enumerate}
\item[\textup{(i)}] 存在一个在 \(\mathcal O_k\) 上平坦的 \(p\)-Lie代数 \(\mathbf L\)，使 \(G\) 同构于 \(\mathscr G_p(\mathbf L)\)。此时，由 \SGARefLink{sga3:VIIB:4.2.2}{4.2.2}，有 \(\mathbf L=\operatorname{Lie}(G)\)。

\item[\textup{(ii)}] 存在投射伪紧 \(k\)-模 \(\omega\)，使 \(G\) 的仿射代数同构于 \(k[[\omega]]\) 对由元素 \(x^p\)，\(x\in\omega\) 生成的闭理想的商。此时，\(\omega\simeq\omega_{G/k}\)。

\item[\textup{\SGAInlineRefTarget{sga3:viib:subsection:4-4:item:iii}(iii)}] 群 \(G\) 的高度至多为 \(1\)，且 \(\omega_{G/k}\) 为投射伪紧 \(k\)-模。

\item[\textup{(iv)}] 群 \(G\) 的高度至多为 \(1\)，且在 \(k\) 上拓扑平坦。
\end{enumerate}}

\noindent\textit{\textup{\SGARefLink{sga3:vii:paragraph:5-1:equation:i}{(i)}} \(\boldsymbol{\Rightarrow}\) \textup{\SGARefLink{sga3:vii:paragraph:5-1:equation:ii}{(ii)}} 的证明。}
置
\[
  \omega=\Gamma^*(\mathbf L)
\]
\textup{（参见 \SGARefLink{sga3:VIIB:1.2.3.D}{1.2.3.D}）}，并令 \(A\) 为 \(k[[\omega]]\) 对由元素 \(x^p\)，\(x\in\omega\) 生成的闭理想的商。每个连续态射
\[
  h:A\longrightarrow C,
  \qquad C\in\operatorname{Alf}/k,
\]
由其在 \(\omega\) 上的限制 \(h'\) 确定，而此限制将 \(\omega\) 映入 \(\sqrt[p]{0_C}\)。因此，从
\[
  \operatorname{Hom}_{\operatorname{Alp}/k}(A,C)
\]
到由 \(\omega\) 至 \(\sqrt[p]{0_C}\) 的连续 \(k\)-线性映射组成的集合
\[
  \operatorname{Hom}_{c}\bigl(\omega,\sqrt[p]{0_C}\bigr)
\]
有一个典范双射。后一个集合典范同构于
\[
  \mathbf L(C)\otimes_C\sqrt[p]{0_C}
\]
\textup{（见 \SGARefLink{sga3:VIIB:3.3:cartier-theorem}{3.3} 的证明）}。因此，蕴含关系 \textup{\SGARefLink{sga3:vii:paragraph:5-1:equation:i}{(i)}} \(\Rightarrow\) \textup{\SGARefLink{sga3:vii:paragraph:5-1:equation:ii}{(ii)}} 来自 \SGARefLink{sga3:VIIA:4.3.1}{4.3.1} 中建立的函子性双射
\[
  \mathbf L(C)\otimes_C\sqrt[p]{0_C}
  \xrightarrow{\ \sim\ }
  \mathscr G_p(\mathbf L)(C)
\]

其余蕴含关系可参看 \SGARefLink{sga3:VIIB:3.3:cartier-theorem}{定理~3.3} 和 VIIA~\SGARefLink{sga3:VIIA:7.4.1}{7.4.1} 中类似的证明。

\medskip
\noindent\textbf{注 4.4.A.}
\phantomsection
\label{sga3:VIIB:4.4.A}
设 \(G\) 为一个无穷小形式 \(k\)-群，仿射代数为 \(A\)，且
\[
  \omega_{G/k}=I_A/I_A^2
\]
为投射伪紧 \(k\)-模。则投影
\[
  I_A\longrightarrow\omega_{G/k}
\]
有一个截面
\[
  \tau:\omega_{G/k}\longrightarrow I_A,
\]
而 \(\tau\) 诱导满的连续代数态射
\[
  \psi:k[[\omega_{G/k}]]\longrightarrow A;
\]
参看 \SGARefLink{sga3:VIIB:3.3:cartier-theorem}{3.3} 中蕴含关系 \textup{\SGARefLink{sga3:viib:subsection:4-4:item:iii}{(iii)}} \(\Rightarrow\) \textup{\SGARefLink{sga3:vii:paragraph:5-1:equation:i}{(i)}} 的证明。\footnote{%
\label{sga3:VIIB:nde:139}编者增加了此注；\SGARefLink{sga3:VIIB:4.4.2}{4.4.2} 中要用到它。}

\subsubsection*{4.4.1 推论}
\phantomsection
\label{sga3:VIIB:4.4.1}

\noindent\textbf{推论 4.4.1.}
{\itshape 若 \(k\) 为特征 \(p>0\) 的域，则函子
\[
  \mathbf L\longmapsto\mathscr G_p(\mathbf L)
\]
给出 \(k\) 上的 \(p\)-Lie代数范畴与高度至多为 \(1\) 的形式 \(k\)-群范畴之间的等价。}

事实上，当 \(G\) 遍历高度至多为 \(1\) 的形式群时，\SGARefLink{sga3:viib:subsection:4-4}{定理~4.4} 将函子
\[
  G\longmapsto\mathscr G_p\bigl(\operatorname{Lie}G\bigr)
\]
与恒等函子等同。同样，\SGARefLink{sga3:VIIB:4.2.2}{4.2.2} 将函子
\[
  \mathbf L\longmapsto
  \operatorname{Lie}\mathscr G_p(\mathbf L)
\]
与恒等函子等同\textup{（另见 VIIA~\SGARefLink{sga3:VIIA:5.5.3}{5.5.3}）}。\footnote{%
\label{sga3:VIIB:nde:140}若 \(k\) 为特征 \(p>0\) 的 Artin 环，同样的证明给出在 \(k\) 上平坦的 \(p\)-Lie代数范畴与高度至多为 \(1\) 且在 \(k\) 上拓扑平坦的形式 \(k\)-群范畴之间的等价。}
